\par\bigskip
\begin{center}
\textbf{VII. Die Jahre 1811, 1812. Gausssche Reihe. Funktionentheorie.}\\[6pt]
\textls[100]{a) Korrespondenz mit F. W. \textsc{Bessel} 1810—1812.}
\end{center}
\noindent\hspace*{\parindent}Unter den Männern, mit denen \textsc{Gauss} Jahrzehnte hindurch in regelmässigem
Briefwechsel gestanden hat, war \textsc{Bessel} derjenige, der seinen mathematischen
Interessen zufolge \textsc{Gauss} am meisten zu bieten und demgemäss
auch von ihm am meisten zu empfangen vermochte. So sind uns namentlich
die Briefe, die zwischen dem 26. August 1810 und dem 5. Mai 1812 zwischen
beiden Männern gewechselt wurden, in hohem Grade wertvoll, da sich \textsc{Gauss}
in ihnen über Gegenstände der Analysis ausgesprochen hat, über die wir
keine andere Äusserung von ihm besitzen.
\par
Angeregt durch das Erscheinen von \textsc{Soldners} Werk über den Integrallogarithmus\footnote{J. \textsc{Soldner}, \textit{Théorie et tables d’une nouvelle fonction transcendante.} München 1809.}
hatte sich \textsc{Bessel} mit dieser schon von \textsc{Euler} und \textsc{Mascheroni}
untersuchten Funktion beschäftigt und berichtet über seine Resultate\footnote{\textit{Briefwechsel Gauss-Bessel}, S. 113—121, Brief vom 26. August 1810.} an\par
\clearpage\SourceLeaf{220}\ArithmeticHead{135}{R}{VIIa. KORRESPONDENZ MIT BESSEL 1810—1812.}
\noindent \textsc{Gauss}, weil dieser einmal den Wunsch geäussert hatte, »die Funktion \(\operatorname{li}x\) für
sehr grosse \(x\) zu kennen, um die schöne Bemerkung des Zusammenhanges
mit den Primzahlen daran prüfen zu können.«
\par
\textsc{Gauss} hatte sich, wie er 1849 an \textsc{Encke} schreibt (Werke II, S. 444),
schon sehr früh (1792—1793) mit dem Problem der Frequenz der Primzahlen
beschäftigt und »erkannte bald, dass unter allen Schwankungen diese Frequenz
durchschnittlich nahe dem Logarithmus verkehrt proportional sei, so dass die
Anzahl aller Primzahlen unter einer gewissen Grenze \(n\) nahe durch das Integral
\begingroup\makeatletter\tagsleft@true\let\veqno\leqno\makeatother
\begin{equation*}
\int\frac{dn}{\log n}\tag{1}
\end{equation*}\endgroup
\noindent ausgedrückt wird.« Nachdem \textsc{Bessel} seine Mitteilungen am 19. Oktober fortgesetzt
hatte, antwortet \textsc{Gauss} mit dem bereits erwähnten Briefe vom 21. Oktober
1810. Eine ausführliche Kritik von \textsc{Bessels} Auffassung des Integrallogarithmus
gibt \textsc{Gauss} erst in dem Briefe vom 5. Mai 1812 (Werke X1,
S. 374—378).
\par
Am 10. März 1811\footnote{\textit{Briefwechsel Gauss-Bessel}, S. 138—143.} berichtet \textsc{Bessel} abermals über eine mathematische
Untersuchung, die sich auf die, wie er sich ausdrückt, »berüchtigten« \textsc{Krampschen}
Fakultäten bezieht. Auf diese Mitteilung erwidert \textsc{Gauss} (6. April 1811,
a. a. O. S. 143), dass er noch nicht Zeit gehabt habe, \textsc{Bessels} Bemerkungen
durchzugehen, und nachdem \textsc{Bessel} (16. Oktober 1811, a. a. O. S. 150) noch
geschrieben, dass er seine Untersuchungen über die Fakultäten ganz umgearbeitet
habe und sie für interessant halte, folgen zwei Briefe von \textsc{Gauss} vom
21. November (Werke X, 1, S. 362) und vom 18. Dezember (ebenda, S. 365) in
denen sich, veranlasst durch kritische Bemerkungen zu \textsc{Bessels} Arbeiten über
Fakultäten und Integrallogarithmus, äusserst wichtige Aussprüche von \textsc{Gauss}
finden.
\par
»Ich arbeite nämlich jetzt«, schreibt \textsc{Gauss} am 21. November, »an einer
Abhandlung für unsere Sozietät, die in etwa sechs Wochen vollendet sein
wird und die Reihe
\begingroup\makeatletter\tagsleft@true\let\veqno\leqno\makeatother
\begin{equation*}
\begin{aligned}
1+\frac{\alpha\mathbin{.}\beta}{1\mathbin{.}\gamma}x
&+\frac{\alpha\mathbin{.}\alpha+1\mathbin{.}\beta\mathbin{.}\beta+1}
{\gamma\mathbin{.}\gamma+1\mathbin{.}1\mathbin{.}2}x^2\\
&+\frac{\alpha\mathbin{.}\alpha+1\mathbin{.}\alpha+2\mathbin{.}\beta\mathbin{.}\beta+1\mathbin{.}\beta+2}
{1\mathbin{.}2\mathbin{.}3\mathbin{.}\gamma\mathbin{.}\gamma+1\mathbin{.}\gamma+2}x^3+\text{etc.}
\end{aligned}\tag{2}
\end{equation*}\endgroup
\noindent betrifft, und auch die Funktionen, die mit \textsc{Kramps} Fakultäten zusammenhängen,
berührt«.
\par
\clearpage\SourceLeaf{221}\ArithmeticHead{136}{L}{L. SCHLESINGER, ÜBER GAUSS' FUNKTIONENTHEORETISCHE ARBEITEN.}
\indent Da hiernach diese beiden Briefe gewissermassen aus dem geistigen Milieu
heraus geschrieben sind, in dem die am 30. Januar 1812 der Sozietät vorgelegte
Abhandlung \textit{Allgemeine Untersuchungen über die Reihe} (2)\footnote{\textit{Disquisitiones generales circa seriem infinitam etc.}, Comm. soc. reg. scient. Gotting. rec. vol. II, 1813, Werke III, S. 123—162; wir zitieren diese Abhandlung in dem ganzen darauf bezüglichen Abschnitt VII kurz als \textit{Disquisitiones}.} redigiert
worden ist, werden wir ihren Inhalt bei der Analyse dieser Abhandlung sorgfältig
zu berücksichtigen haben. Ehe wir nun auf diese Analyse eingehen,
berichten wir über die Vorgeschichte der \textit{Disquisitiones}.
